\linespread{1.4}
\begin{center}
    {\LARGE \textbf{Deep Learn} \hspace{0.5cm}}\\
	\vspace{0.5cm}
	\textbf{Pierre Augusto Ré Martho}\\
	Prof: Nina Sumiko Tomita Hirata\\
	2025\\
	\vspace{0.15cm}
	\textbf{Atividade 2}
\end{center}    
\rule{\textwidth}{0.5pt}


%%%%%%%%%%%%%%%%%%%%%%%%%%%%%%%%%%%%%%%%%%%%%%%%%%%%%%%%%%%%%%%%%%%%%%%%%%%%%%%%
\noindent O notebook com os códigos em Python utilizados para a exploração de dados em anexo: \texttt{Atividade2.ipynb}.

\section*{Exploração de Dados}
\subsection*{Modelos Adotados}
\noindent\textbf{Dataset: MNIST} \\
O conjunto de dados MNIST foi utilizado, conforme documentação oficial: \url{https://docs.pytorch.org/vision/main/datasets.html}

\begin{figure}[H]
    \centering
    \includegraphics[width=0.3\textwidth]{output1.png}
    \caption{Exemplos de imagens do dataset MNIST.}
    \label{fig:mnist_examples}
\end{figure}


\noindent O download e carregamento do dataset foi realizado com o seguinte comando:
\begin{verbatim}
datasets.MNIST(root="dataset/", download=True, train=True, transform=transforms.ToTensor())
\end{verbatim}

\noindent Referências utilizadas para implementação e exemplos:
\begin{itemize}
    \item \url{https://www.datacamp.com/tutorial/pytorch-cnn-tutorial}
    \item \url{https://docs.pytorch.org/tutorials/beginner/blitz/cifar10_tutorial.html}
    \item \url{https://nextjournal.com/gkoehler/pytorch-mnist}
\end{itemize}

\noindent\textbf{Modelos Utilizados:}
\begin{itemize}
    \item \textbf{Rede Neural Simples}: Composta por 3 camadas totalmente conectadas. A entrada possui 784 neurônios (28x28 pixels) e a saída 10 neurônios, representando as classes dos dígitos.
    \item \textbf{Rede Neural Convolucional (CNN)}: Composta por 2 camadas convolucionais seguidas de camadas de pooling e uma camada linear de saída para classificação.
\end{itemize}

\noindent\textbf{Otimizador}: \texttt{optim.Adam} com \texttt{learning rate = 0.001} \\
\noindent\textbf{Função de perda}: \texttt{nn.CrossEntropyLoss()} \\
\noindent\textbf{Número de épocas}: 10 \\
\noindent\textbf{Batch size}: 100 \\

\noindent Foram utilizados subconjuntos do MNIST com diferentes quantidades de objetos por classe: [100, 300, 600, 900].

\noindent Os resultados a seguir mostram a evolução da função de perda (loss) para ambos os modelos ao longo das 10 épocas, considerando diferentes quantidades de exemplos por classe.

% Espaço reservado para imagens dos resultados
\begin{figure}[H]
    \centering
    \includegraphics[width=1.\textwidth]{output2.png}
    \caption{Evolução da função de perda (loss) ao longo das épocas para diferentes quantidades de exemplos por classe nos modelos MLP e CNN.}
    \label{fig:loss_evolution}
\end{figure}

A função de perda diminui para ambos os modelos, sendo ligeiramente melhor para o modelo CNN. Além disso, quanto maior o número de exemplos por classe, mais rapidamente o modelo atinge valores menores de loss, indicando melhor desempenho já nas primeiras épocas do treinamento.
Até a 10ª época, a loss já convergiu para um valor suficientemente baixo para as classes, embora continue decaindo lentamente. Com um maior número de exemplos por classe, o desempenho foi superior.

Aplicando o modelo a um conjunto de validação e observando o desempenho ao longo das épocas nos conjuntos de treinamento e teste, temos a Figura \ref{fig:loss_evolution_test}.


\begin{figure}[h]
    \centering
    \includegraphics[width=1.\textwidth]{output3.png}
    \caption{Evolução da função de perda (loss) ao longo das épocas para o conjunto de teste, utilizando os modelos treinados.}
    \label{fig:loss_evolution_test}
\end{figure}


Para a NN, a loss no conjunto de teste começa a aumentar a partir da época 5, indicando o início de overfitting. Isso é esperado, pois o modelo utiliza apenas camadas lineares sobre o vetor de pixels da imagem, limitando a generalização, mesmo para imagens relativamente simples como as do MNIST. Esse é o início da "boca de jacaré".

No caso do modelo convolucional (CNN), a função de perda no conjunto de teste segue melhor a de treinamento até aproximadamente a época 10, apresentando apenas um leve aumento posterior. Isso demonstra que a CNN generaliza melhor. Para conjuntos maiores (por exemplo, 900 exemplos por classe), treinar por cerca de 5 épocas já é suficiente para obter bom desempenho, visto que o MNIST é um conjunto relativamente simples.

\vspace{0.5cm}

\section*{Explorações adicionais}

\noindent a) Imagens modificadas

\noindent Para testar a robustez dos modelos, foi adicionado ruído "sal e pimenta" às imagens do MNIST, alterando aleatoriamente pixels para preto ou branco. Assim, avaliou-se o desempenho dos modelos com dados corrompidos.

Exmplos imgs modificadas

\begin{figure}[H]
    \centering
    \includegraphics[width=0.3\textwidth]{output4.png}
    \caption{Exemplos de imagens do dataset MNIST modificadas.}
    \label{fig:mnist_examples}
\end{figure}

Repetindo o treinamento com a NN e com a CNN para diferentes tamanhos de amostra, e aplicando-as ao conjunto de teste, observa-se a evolução da loss abaixo:

\begin{figure}[h]
    \centering
    \includegraphics[width=1.\textwidth]{loss_evolution_test_2.png}
    \caption{Evolução da função de perda (loss) ao longo das épocas para o conjunto de teste, utilizando os modelos treinados para as imagems modificadas.}
    \label{fig:loss_evolution_test_2}
\end{figure}

A loss converge bem mais devagar do que no caso simples. Em contrapartidas, a evolução da loss na amostra de treino mostrou pouco sinal de overfitting. Como cada uma das imagens foi alterada de forma diferente, o modelo tende a ter mais dificuldade para memorizar os dados, o que pode explicar a ausência de overfitting e uma melhor generalização, mesmo no caso da NN.

